\chapter{Mammography}
\index{Mammography}

\section*{Definition and Overview}
Mammography is a specialised medical imaging technique that uses low-dose X-rays to visualise the internal structures of the breast. It is the primary screening tool for the early detection of breast cancer, capable of identifying tumours, calcifications, and architectural distortions long before they can be felt on physical examination. Mammography is broadly classified into screening mammography (performed on asymptomatic individuals to detect early signs of breast cancer) and diagnostic mammography (performed to investigate specific breast symptoms, such as a lump, pain, skin dimpling, or nipple discharge, or to clarify abnormalities found on a screening mammogram). Modern mammography utilises digital technology (digital mammography) and often includes digital breast tomosynthesis (three-dimensional mammography), which captures multiple X-ray images from different angles to construct a detailed 3D reconstruction of the breast tissue.

\section*{Epidemiology}
Breast cancer is the most common cancer diagnosed in women globally and in many countries, representing a major public health concern. In many countries, approximately 1 in 7 women will be diagnosed with breast cancer in their lifetime. Screening mammography is widely implemented through national programmes like the local breast-screening programme, which targets women aged 50--74 years for free, two-yearly screening, although women aged 40--49 and 75 and over are also eligible. These screening programmes have significantly contributed to a reduction in breast cancer mortality by enabling early detection, when treatments are most effective and less invasive. While breast cancer predominantly affects women, it can also occur in men (representing less than 1\% of cases), in whom mammography is used diagnostically to evaluate breast masses.

\section*{Aetiology and Pathophysiology}
Mammography is indicated to detect breast cancer, which arises from the malignant transformation of epithelial cells lining the ducts (ductal carcinoma) or lobules (lobular carcinoma) of the breast. The X-ray imaging depends on tissue density differentials within the mammary gland.
\begin{itemize}
    \item \textbf{Tissue density differentials:} Normal breast tissue consists of fat, fibrous connective tissue, and glandular tissue. Fat is radiolucent (appears dark on the mammogram), whereas glandular and fibrous tissue are radiopaque (appear light). Breast cancer cells and microcalcifications are also radiopaque, making them visible against a dark background of fat.
    \item \textbf{Dense breast tissue limitation:} In women with dense breasts (high proportion of glandular and fibrous tissue), both normal tissue and tumours appear white, which can mask underlying cancers and reduce mammographic sensitivity.
    \item \textbf{Microcalcifications:} Tiny calcium deposits within the breast ducts that appear as bright white specks; specific patterns (clustered, pleomorphic) can indicate ductal carcinoma in situ or invasive cancer.
    \item \textbf{Architectural distortion:} Disruption of the normal radial pattern of the breast tissue, which can indicate an invasive carcinoma even in the absence of a distinct mass.
\end{itemize}

\section*{Clinical Features}
Mammography itself is an investigative procedure, but its indications and findings correlate with specific clinical presentations.
\begin{itemize}
    \item Screening mammography is performed in asymptomatic individuals, with no physical breast changes, skin abnormalities, or pain.
    \item Diagnostic mammography is indicated for clinical symptoms such as a palpable breast lump or focal breast thickness.
    \item Skin changes including localised dimpling, puckering, or retraction (resembling an orange peel, known as peau d'orange) over the breast.
    \item Nipple abnormalities, including recent nipple inversion, scaling, redness, or spontaneous, bloody nipple discharge.
    \item Mammographic features of breast cancer include an irregular, spiculed mass (a white shape with fuzzy, radiating borders), focal asymmetry, or clustered microcalcifications.
    \item Post-procedure discomfort is common, characterised by mild, transient breast soreness or bruising due to the mechanical compression required during the scan.
\end{itemize}

\section*{Differential Diagnosis}
On mammographic imaging, malignant findings must be differentiated from benign breast lesions:
\begin{itemize}
    \item \textbf{Fibroadenoma:} A common, benign, solid breast tumour that typically appears on mammography as a well-circumscribed, oval mass, sometimes with classic popcorn-like calcifications in older women.
    \item \textbf{Simple breast cyst:} A fluid-filled sac that appears as a round or oval mass; distinguished from solid masses by breast ultrasound rather than mammography.
    \item \textbf{Fat necrosis:} A benign condition resulting from trauma or surgery that can mimic cancer by presenting as a firm mass with irregular calcifications or architectural distortion.
    \item \textbf{Radial scar:} A benign breast lesion characterised by a star-like pattern of architectural distortion, which is mammographically indistinguishable from invasive breast cancer and requires surgical biopsy.
    \item \textbf{Mastitis or abscess:} Inflammatory conditions that cause generalised tissue density and skin thickening, mimicking inflammatory breast cancer.
\end{itemize}

\section*{When to Seek Medical Care}
Women should participate in regular screening mammography according to guidelines. Any individual who notices new or unusual breast changes should consult their general practitioner immediately for a diagnostic evaluation.

\textbf{Contact your local emergency services for:}
\begin{itemize}
    \item Severe, acute chest pain or shortness of breath (unrelated to mammography, but representing urgent medical emergencies).
    \item Signs of severe systemic infection or sepsis (such as high fever, confusion, rapid heart rate) in a patient with an infected breast abscess.
\end{itemize}

Other non-emergency reasons to see a doctor include:
\begin{itemize}
    \item Finding a new lump or area of thickening in the breast or underarm.
    \item Noticing skin dimpling, puckering, or redness on the breast.
    \item Recent changes in the nipple, such as retraction, inversion, scaling, or discharge (especially if bloody or spontaneous).
    \item Persistent breast pain that is not related to the menstrual cycle.
    \item Receiving a mammogram result letter indicating an abnormality that requires further diagnostic testing.
\end{itemize}

\section*{Diagnosis and Investigations}
Mammography is a key component of the triple test used to diagnose breast anomalies.
\begin{itemize}
    \item \textbf{Triple test approach:} The gold standard for investigating breast symptoms, combining clinical breast examination, breast imaging (mammography and/or ultrasound), and tissue biopsy.
    \item \textbf{Mammographic views:} Standard screening includes two views of each breast: the craniocaudal (top-down) view and the mediolateral oblique (angled side-to-side) view.
    \item \textbf{Digital breast tomosynthesis:} Takes multiple low-dose X-rays at various angles to create a 3D image, improving detection rates and reducing false-positive recalls, especially in dense breast tissue.
    \item \textbf{Breast ultrasound:} Utilised as an adjunct to mammography to differentiate solid masses from fluid-filled cysts, particularly in younger women (under 35 years) who have naturally dense breasts.
    \item \textbf{Breast MRI:} Reserved for high-risk patients, screening women with \textit{BRCA} gene mutations, or evaluating breast implant integrity and disease extent.
    \item \textbf{BI-RADS reporting system:} The Breast Imaging-Reporting and Data System, which standardises mammogram results from category 0 (incomplete, needs further imaging) to category 6 (known biopsy-proven malignancy).
\end{itemize}

\section*{Management}
The management related to mammography includes scheduling regular screening, interpreting results, and pursuing follow-up diagnostics or cancer treatment.
\begin{itemize}
    \item \textbf{Screening participation:} Eligible women are encouraged to book free screening mammograms every two years through BreastScreen the affected region or local equivalent programmes.
    \item \textbf{Recall assessment:} If a screening mammogram shows an abnormality, the patient is recalled to a diagnostic clinic for further views, ultrasound, clinical exam, and biopsy if necessary.
    \item \textbf{Fine needle aspiration or core biopsy:} Performed under ultrasound or stereotactic (mammographic) guidance to obtain cells or tissue cores for pathological analysis.
    \item \textbf{Surgical consultation:} Indicated if biopsy confirms malignancy or an atypical lesion (such as atypical ductal hyperplasia) that requires surgical excision.
    \item \textbf{Oncological management:} If breast cancer is diagnosed, treatment is coordinated by a multidisciplinary team and may include lumpectomy, mastectomy, chemotherapy, hormone therapy, targeted therapy, and radiotherapy.
    \item \textbf{High-risk surveillance:} Personalised screening plans, including annual mammograms and MRIs, for women with a strong family history or proven genetic mutations.
\end{itemize}

\section*{Complications}
\begin{itemize}
    \item \textbf{False positives and anxiety:} Mammography can identify suspicious areas that turn out to be benign upon further imaging or biopsy, causing significant psychological distress.
    \item \textbf{False negatives:} Cancers can be missed on mammography, particularly in dense breasts, leading to delayed diagnosis and false reassurance.
    \item \textbf{Overdiagnosis and overtreatment:} Detection of slow-growing, non-progressive cancers (such as some cases of ductal carcinoma in situ) that might never have caused symptoms or harmed the patient in their lifetime.
    \item \textbf{Radiation exposure:} Mammography involves exposure to ionising radiation; however, the dose is extremely low (equivalent to about 7 weeks of natural background radiation) and the benefits of early detection far outweigh the risks.
    \item \textbf{Mechanical discomfort:} The physical compression of the breast between plates can cause brief, moderate pain and, rarely, minor skin bruising.
\end{itemize}

\section*{Prognosis and Outcomes}
Mammography screening has been shown to reduce breast cancer mortality by 20--30\% in the target age group of 50--74 years. Early detection through mammography is associated with a highly favourable prognosis; the 5-year survival rate for localised breast cancer (stage I) is close to 100\%, compared to approximately 30\% for metastatic (stage IV) disease. Patients diagnosed early also have more treatment options, such as breast-conserving surgery (lumpectomy) instead of mastectomy, and may avoid the need for chemotherapy.

\section*{Prevention and Self-Care}
\begin{itemize}
    \item Attend routine mammograms every two years if you are in the target age group (50 to 74 years), or discuss an appropriate screening schedule with your doctor if you are younger or have risk factors.
    \item Practice breast self-awareness: get to know the normal look and feel of your breasts and report any changes to your general practitioner immediately.
    \item Schedule your mammogram for the week after your menstrual period, when breasts are typically less tender, to minimise compression discomfort.
    \item Do not apply deodorant, perfume, powder, or lotion under your arms or on your breasts on the day of the mammogram, as these products can contain metallic particles that appear as calcifications on the X-ray.
    \item Wear a two-piece outfit (such as a top and trousers) so that you only need to remove your top during the procedure.
    \item Discuss your family history of breast and ovarian cancer with your doctor to determine if you need early or enhanced screening.
\end{itemize}

\section*{Sources and Further Reading}
The following organisations provide authoritative, evidence-based information on mammography and breast cancer screening:
\begin{itemize}
    \item BreastScreen Australia. \textit{National breast cancer screening program guidelines.} https://www.health.gov.au/our-work/breastscreen-australia-program
    \item Cancer Council Australia. \textit{Early detection and screening for breast cancer.} https://www.cancer.org.au/
    \item NHS. \textit{Breast cancer screening (mammogram) guidelines.} https://www.nhs.uk/conditions/breast-screening-mammogram/
    \item Mayo Clinic. \textit{Mammogram: overview, expectations, and results.} https://www.mayoclinic.org/tests-procedures/mammogram/
    \item American Cancer Society. \textit{Mammograms and breast cancer screening.} https://www.cancer.org/cancer/types/breast-cancer/screening-tests-and-early-detection/mammograms.html
\end{itemize}
